\subsection{Attachment of a topological space}

Let X, B be topological spaces, let A be a subspace of B and let \(f: A \rightarrow X\) be a continuous map. Let Y denote the topological sum of the spaces \(Y_{1} = X\) and \(Y_{2} = B\) and identify \(Y_{1}\) and \(Y_{2}\) with subspaces of Y by the canonical injections. Let R be the equivalence relation in Y which is the finest for which the elements a of \(Y_{2}\) and \(f(a)\) of \(Y_{1}\) are equivalent, for every \(a \in A\) . The relation R is the equality relation in \(Y_{1}\) . Let \(x \in Y_{1}\) and let \(b \in Y_{2}\) ; we have x R b if and only if \(b \in A\) and \(f(b) = x\) . Let \(b, b' \in Y_{2}\) ; in order that one have b R \(b'\) , it is necessary and sufficient that \(b = b'\) , or that \(b, b' \in A\) and \(f(b) = f(b')\) .

DEFINITION 9. --- The quotient topological space Y/R is called the space obtained by attaching to the space X the space B along the map f. It is denoted \(X \cup_{f} B\) .

Denote by \(\alpha_{f}\colon\mathrm{X}\to\mathrm{X}\cup_{f}\mathrm{B}\) and \(\beta_{f}\colon\mathrm{B}\to\mathrm{X}\cup_{f}\mathrm{B}\) the restrictions to \(\mathrm{Y}_{1}\) and \(\mathrm{Y}_{2}\) of the canonical surjection of Y onto Y/R. We have \(\alpha_{f}\circ f=\beta_{f}\mid\mathrm{A}\) .

Remarks. --- 1) The map \(\alpha_{f}\) is injective. For every subset U of X, we have \(\bar{\alpha}_{f}^{-1}(\alpha_{f}(U)) = U\) and \(\bar{\beta}_{f}^{-1}(\alpha_{f}(U)) = f^{-1}(U)\) . If A is an open (resp. closed) subset of B, these relations imply that \(\alpha_{f}\) is an open (resp. closed) map from X into \(X \cup_{f} B\) .

Let U be an open subset of X and let V be an open set of B such that \(f^{-1}(U) = V \cap A\) . The union of the subsets U of \(Y_{1}\) and V of \(Y_{2}\) is an open saturated subset of Y; its image in \(X \cup_{f} B\) is therefore open and its trace on \(\alpha_{f}(X)\) is \(\alpha_{f}(U)\) . Consequently, the map \(\alpha_{f}\) defines a homeomorphism of X onto its image in \(X \cup_{f} B\) .

\begin{enumerate}
\def\labelenumi{\arabic{enumi})}
\setcounter{enumi}{1}
\tightlist
\item
  Let V be a subset of B. We have \(\bar{\alpha}_{f}^{-1}(\beta_{f}(V)) = f(V \cap A)\) and \(\bar{\beta}_{f}^{-1}(\beta_{f}(V)) = V \cup \bar{f}(f(V \cap A))\) . If A is closed in B and if the map f is closed, the map \(\beta_{f}\) is therefore closed. Similarly, if A is open in B and if the map f is open, the map \(\beta_{f}\) is open.
\end{enumerate}

In these two cases, the map \(\beta_{f}\) then induces a homeomorphism of \(B\setminus A\) onto its image.

\begin{enumerate}
\def\labelenumi{\arabic{enumi})}
\setcounter{enumi}{2}
\tightlist
\item
  Suppose that A is closed and that the map \(f\) is proper. We have just seen that the map \(\beta_f\) is closed. For every \(x \in X\) , \(\overline{\beta_f}(\alpha_f(x)) = \overline{f}(x)\) . It is therefore a quasi-compact subset of A (TG, I, p. 75, theorem 1), hence also of B. For every \(b \in B\) , \(\overline{\beta_f}(\beta_f(b))\) equal to \(\{b\}\) if \(b \in B - A\) , and \(\overline{f}(f(b))\) if \(b \in A\) ; in both cases, it is a quasi-compact subset of B. The fibers of the map \(\beta_{f}\) are therefore quasi-compact; consequently (loc. cit.), the map \(\beta_{f}\) is proper.
\end{enumerate}

The map \(\alpha_{f}\) is also proper (TG, I, p. 72, prop. 2). It follows that the canonical map of Y onto X \(\cup_{f}\) B is proper.

Examples. --- 1) Let X and Y be topological spaces and \(f\colon X\to Y\) a continuous map. The cylinder \(\operatorname{Cyl}(f)\) is none other than the space obtained by attaching to the space Y the space \(X\times I\) along the map \(f\circ \mathrm{pr}_1\) of \(X\times \{1\}\) in Y.

\begin{enumerate}
\def\labelenumi{\arabic{enumi})}
\setcounter{enumi}{1}
\tightlist
\item
  Let X be a topological space, let B be a topological space and let A be a subspace of B, let \(f: A \to X\) be a continuous map which induces a homeomorphism of A onto its image.
\end{enumerate}

Let us set \(\mathrm{A}' = f(\mathrm{A})\); let \(f'\) be the map from \(\mathrm{A}'\) into B which associates to every \(x \in \mathrm{A}'\) the unique preimage of \(x\) under \(f\). The map \(f' \circ f\) is continuous; since \(f\) is strict, the map \(f'\) is continuous. There exists a unique continuous map \(u\) from the space \(\mathrm{X} \cup_f \mathrm{B}\) to the space \(\mathrm{B} \cup_{f'} \mathrm{X}\) which maps \(\alpha_f(x)\) onto \(\beta_{f'}(x)\) for \(x \in \mathrm{X}\) and \(\beta_f(b)\) onto \(\alpha_{f'}(b)\) for \(b \in \mathrm{B}\); it is a homeomorphism which maps the subspace \(\alpha_f(\mathrm{X})\) onto the subspace \(\beta_{f'}(\mathrm{X})\).

PROPOSITION 8. --- Let Z be a topological space.

\begin{enumerate}
\def\labelenumi{\alph{enumi})}
\item
  Let \(u \colon X \to Z\) and \(v \colon B \to Z\) be continuous maps such that \(v(a) = u(f(a))\) for all \(a \in A\). There then exists a unique continuous map \(w \colon X \cup_f B \to Z\) such that \(w \circ \alpha_f = u\) and \(w \circ \beta_f = v\).
\item
  Let \(\sigma\colon\mathbf{X}\times\mathbf{I}\to\mathbf{Z}\) and \(\tau\colon\mathbf{B}\times\mathbf{I}\to\mathbf{Z}\) be homotopies. It is assumed that \(\tau\) is fixed on A and that \(\sigma(f(a),t)=\tau(a,t)\) for all \(a\in\mathbf{A}\) and every \(t\in\mathbf{I}\). There then exists a unique homotopy \(\eta\colon(\mathbf{X}\cup_{f}\mathbf{B})\times\mathbf{I}\to\mathbf{Z}\) such that \(\eta(\alpha_{f}(x),t)=\sigma(x,t)\) and \(\eta(\beta_{f}(b),t)=\tau(b,t)\) for \(x\in\mathbf{X}\), \(b\in\mathbf{B}\), and \(t\in\mathbf{I}\). This homotopy is fixed on \(\beta_{f}(\mathbf{A})\).
\end{enumerate}

The proposition follows immediately from the definition of a quotient space and from Prop. 10 (I, p. 19).

COROLLARY 1. --- Let \(B'\) be a topological space, let \(A'\) be a subspace of \(B'\), and let \(f': A' \to X\) be a continuous map. Let \(v: B \to B'\) be a continuous map such that \(v(A) \subset A'\) and \(f = f' \circ (v \mid A)\). Let \(w: X \cup_f B \to X \cup_{f'} B'\) be the unique continuous map such that \(w \circ \alpha_f = \alpha_{f'}\) and \(w \circ \beta_f = \beta_{f'} \circ v\). If \(v\) defines a homotopy equivalence of the pair \((B, A)\) onto the pair \((B', A')\), then \(w\) defines a homotopy equivalence of the pair \((X \cup_f B, \alpha_f(X))\) onto the pair \((X \cup_{f'} B', \alpha_{f'}(X))\).

Let \(v' \colon \mathrm{B}' \to \mathrm{B}\) be a continuous map, let \(\tau \colon \mathrm{B} \times \mathbf{I} \to \mathrm{B}\) be a homotopy fixed on A connecting \(\mathrm{Id}_{\mathrm{B}}\) to \(v' \circ v\), and let \(\tau' \colon \mathrm{B}' \times \mathbf{I} \to \mathrm{B}'\) be a homotopy fixed on \(\mathrm{A}'\) connecting \(\mathrm{Id}_{\mathrm{B}'}\) to \(v \circ v'\). For \(a' \in \mathrm{A}'\) and \(a = v'(a')\), we have \(a' = v(a)\) and \(\beta_f(v'(a')) = \beta_f(a) = \alpha_f(f(a)) = \alpha_f(f'(a'))\). By Proposition 8(a), there exists a unique map \(w' \colon X \cup_{f'} B' \to X \cup_f B\) such that \(w' \circ \alpha_{f'} = \alpha_{f'}\) and \(w' \circ \beta_{f'} = \beta_f \circ v'\). By Proposition 8(b), there exists a unique homotopy \(\eta \colon (\mathrm{X} \cup_f \mathrm{B}) \times \mathbf{I} \to X \cup_{f'} B'\) such that \(\eta(\alpha_f(x), t) = \alpha_{f'}(x)\) and \(\eta(\beta_f(b), t) = \beta_{f'}(\tau(b, t))\) for \(x \in X\), \(b \in B\), and \(t \in I\). There likewise exists a unique homotopy \(\eta' \colon (X \cup_{f'} B') \times I \to X \cup_f B\) such that \(\eta'(\alpha_{f'}(x), t) = \alpha_f(x)\) and \(\eta'(\beta_{f'}(b), t) = \beta_f(\tau'(b, t))\) for \(x \in X\), \(b \in B'\), and \(t \in I\). The map from \((X \cup_f B) \times I\) into \(X \cup_f B\) given by \((x, t) \mapsto \eta'(\eta(x, t), t)\) is then a homotopy fixed on \(\alpha_f(X)\) connecting the identity map of \(X \cup_f B\) to the map \(w' \circ w\). Similarly, the map from \((X \cup_{f'} B') \times I\) into \(X \cup_{f'} B'\) given by \((x, t) \mapsto \eta(\eta'(x, t), t)\) is a homotopy fixed on \(\alpha_{f'}(X)\) connecting the identity map of \(X \cup_{f'} B'\) to the map \(w \circ w'\). The corollary follows from this.

Example 3. --- Let i denote the canonical injection of A into B and take as space B' the mapping cylinder of i; denote \(\alpha_{i}\colon\mathrm{A}\times\mathbf{I}\to\mathrm{Cyl}(i)\) and \(\beta_{i}\colon\mathrm{B}\to\mathrm{Cyl}(i)\) the canonical maps and \(r_{i}\colon\mathrm{Cyl}(i)\to\mathrm{B}\) the canonical retraction of the cylinder \(\mathrm{Cyl}(i)\) on its base. Let \(A_{0}\) be the subspace \(\alpha_{i}(\mathrm{A}\times\{0\})\) of \(\mathrm{Cyl}(i)\) and denote by \(f_{0}\colon\mathrm{A}_{0}\to\mathrm{X}\) the map \(f\circ r_{i}\). For \(a\in A\), we have

\[
\beta_ {f} \circ r _ {i} (\alpha_ {i} (a, 0)) = \beta_ {f} (a) = \alpha_ {f} (f _ {0} (\alpha_ {i} (a, 0))) = \alpha_ {f} (f (a)).
\]

Let \(\eta\colon X\cup_{f_0}\mathrm{Cyl}(i)\to X\cup_fB\) be the unique continuous map such that \(\eta\circ\alpha_{f_0}=\alpha_f\) and \(\eta\circ\beta_{f_0}=\beta_f\circ r_i\). Suppose the pair (B, A) has the homotopy extension property. Then the map \(r_i\) defines a homotopy equivalence of the pair \((\mathrm{Cyl}(i),\mathrm{A}_0)\) onto the pair (B,A) (III, p. 243, th. 1). By Corollary 1, the map \(\eta\) then defines a homotopy equivalence of the pair \((X\cup_{f_0}\mathrm{Cyl}(i),\alpha_{f_0}(X))\) onto the pair \((X\cup_fB,\alpha_f(X))\). In particular, \(\eta\) is a homotopy equivalence.

COROLLARY 2. --- Let \(X'\) be a topological space, let \(u: X \to X'\) be a continuous map, and set \(f' = u \circ f\). Let \(w: X \cup_f B \to X' \cup_{f'} B\) be the unique continuous map such that \(w \circ \alpha_f = \alpha_{f'} \circ u\) and \(w \circ \beta_f = \beta_{f'}\). If \(u\) defines a homotopy equivalence of the pair \((X, f(A))\) onto the pair \((X', f'(A))\), then \(w\) defines a homotopy equivalence of the pair \((\mathrm{X}\cup_f\mathrm{B},\beta_f(\mathrm{B}))\) onto the pair \((\mathrm{X}'\cup_{f'}\mathrm{B},\beta_{f'}(\mathrm{B}))\).

The proof is analogous to that of Corollary 1.

PROPOSITION 9. --- Let X and B be topological spaces, let A be a subspace of B, and let f: A → X be a continuous map.

\begin{enumerate}
\def\labelenumi{\alph{enumi})}
\item
  If the pair (B, A) has the homotopy extension property, then so does the pair \((\mathrm{X} \cup_{f} \mathrm{B}, \alpha_{f}(\mathrm{X}))\) .
\item
  Suppose that the map f is injective and strict. If the pair \((\mathrm{X}, f(\mathrm{A}))\) has the homotopy extension property, so does the pair \((\mathrm{X} \cup_{f} \mathrm{B}, \beta_{f}(\mathrm{B}))\) .
\item
  Suppose the pair (B, A) has the homotopy extension property. Let Z be a topological space, let \(u: X \cup_{f} B \to Z\) be a continuous map and let \(\sigma: \alpha_{f}(X) \times I \to Z\) be a homotopy whose term is the restriction of u to the subspace \(\alpha_{f}(X)\).
\end{enumerate}

Let us set \(v_{1} = u \circ \alpha_{f}\) and denote by \(\tau_{1} \colon \mathbf{X} \times \mathbf{I} \to \mathbf{Z}\) the map given by \((x,t) \mapsto \sigma(\alpha_{f}(x),t)\). Set \(v_{2} = u \circ \beta_{f}\). Since \(\beta_{f} \mid \mathrm{A} = \alpha_{f} \circ f\), the map from \(\mathrm{A} \times \mathbf{I}\) into \(\mathrm{Z}\) which maps \((a,t)\) onto \(\sigma(\alpha_{f}(f(a)),t)\) for \(a \in \mathrm{A}\) and \(t \in \mathbf{I}\) is a homotopy whose term is equal to the map \(v_{2} \mid \mathrm{A}\). Since the pair \((\mathrm{B},\mathrm{A})\) has the homotopy extension property, there exists a homotopy \(\tau_{2} \colon \mathbf{B} \times \mathbf{I} \to \mathbf{Z}\) whose term is the map \(v_{2}\) and such that \(\tau_{2}(a,t) = \sigma(\alpha_{f}(f(a)),t)\) for \((a,t) \in \mathrm{A} \times \mathbf{I}\).

For \(a \in \mathbf{A}\) and \(t \in \mathbf{I}\), we have \(\tau_2(a, t) = \tau_1(f(a), t)\). By Proposition 8, there exists a unique continuous map \(\sigma: (\mathrm{X} \cup_f \mathrm{B}) \times \mathbf{I} \to \mathbf{Z}\) such that \(\sigma(\alpha_f(x), t) = \sigma_1(x, t)\) and \(\sigma(\beta_f(b), t) = \sigma_2(b, t)\), for \(x \in \mathrm{X}\), \(b \in \mathrm{B}\), and \(t \in \mathbf{I}\). This is a homotopy with term \(u\) which extends \(\tau\), whence assertion a).

\begin{enumerate}
\def\labelenumi{\alph{enumi})}
\setcounter{enumi}{1}
\tightlist
\item
  In view of Example 2 (III, p. 249), assertion b) follows from assertion a).
\end{enumerate}

